\chapter{행렬표현과 기저변환, 유사}

\chapterintro{같은 선형사상도 기저를 어떻게 잡는지에 따라 전혀 다른 숫자 배열로 보입니다. 그러나 표현이 달라져도 사상 자체의 구조는 변하지 않습니다. 이 장에서는 행렬표현(matrix representation), 기저변환(change of basis), 유사(similarity)를 하나의 흐름으로 묶어, ``표현의 변화''와 ``본질의 보존''을 동시에 정리하겠습니다.}
\chaptergoals{선형사상의 행렬표현을 좌표식으로 엄밀하게 유도할 수 있다.}{기저변환행렬을 이용해 유사공식 $A'=P^{-1}AP$를 증명하고 계산할 수 있다.}{유사 불변량을 사용해 행렬의 동형적 구조를 판정하는 기초를 세울 수 있다.}

이 장에서는 기본 체를 $\mathbb{F}$로 두고, 모든 벡터공간은 유한차원이라 가정하겠습니다.

\sectiontemplate{선형사상의 행렬표현(matrix representation)}{기저를 고정하면 선형사상 계산은 좌표벡터의 행렬곱으로 바뀝니다. 이 절에서는 그 변환 규칙을 정확히 세우고, 합성사상과 행렬곱이 왜 같은 연산인지 확인하겠습니다.}{\item 좌표벡터와 표현행렬의 정의를 정확히 쓸 수 있습니다.\item 좌표식 $[T(x)]_{\mathcal C}= [T]_{\mathcal C\leftarrow\mathcal B}[x]_{\mathcal B}$를 완전하게 증명할 수 있습니다.\item 합성사상과 행렬곱의 대응을 계산으로 확인할 수 있습니다.}{\item 기저와 좌표표현\item 선형사상의 합성}

\begin{definition}
$V$의 순서기저(ordered basis) $\mathcal{B}=(v_1,\dots,v_n)$가 주어졌을 때, 임의의 $x\in V$를
$$
x=\sum_{j=1}^n \alpha_j v_j
$$
로 유일하게 쓰면, 열벡터
$$
[x]_{\mathcal B}=
\begin{bmatrix}
\alpha_1\\
\vdots\\
\alpha_n
\end{bmatrix}
$$
를 $x$의 $\mathcal{B}$-좌표벡터라 한다.
\end{definition}

\begin{definition}
$T:V\to W$가 선형사상이고 $\mathcal{B}=(v_1,\dots,v_n)$, $\mathcal{C}=(w_1,\dots,w_m)$가 각각 $V,W$의 순서기저일 때,
$$
A=[T]_{\mathcal C\leftarrow\mathcal B}\in M_{m\times n}(\mathbb F)
$$
를 다음과 같이 정한다. 각 $j$에 대해 $T(v_j)$의 $\mathcal C$-좌표벡터를 $A$의 $j$번째 열로 둔다. 즉
$$
A=\begin{bmatrix}[T(v_1)]_{\mathcal C}&\cdots&[T(v_n)]_{\mathcal C}\end{bmatrix}.
$$
이 $A$를 $T$의 표현행렬이라 한다.
\end{definition}

\begin{theorem}[좌표식]
$T:V\to W$가 선형이고 위 표기에서 $A=[T]_{\mathcal C\leftarrow\mathcal B}$라 하면, 모든 $x\in V$에 대해
$$
[T(x)]_{\mathcal C}=A[x]_{\mathcal B}
$$
가 성립한다.
\end{theorem}

\paragraph{증명 전략.}
$x$를 $\mathcal B$-기저로 전개한 뒤 선형성을 적용합니다. 그리고 각 $T(v_j)$를 $\mathcal C$-기저로 다시 전개하여 계수를 정리하면 행렬곱의 성분식이 그대로 나타납니다.

\begin{proof}
$A=(a_{ij})$라 두자. 정의에 의해 각 $j$에 대해
$$
T(v_j)=\sum_{i=1}^m a_{ij}w_i
$$
이다. 임의의 $x\in V$를
$$
x=\sum_{j=1}^n x_j v_j
$$
로 쓰면 선형성에 의해
$$
T(x)=\sum_{j=1}^n x_jT(v_j)
=\sum_{j=1}^n x_j\sum_{i=1}^m a_{ij}w_i
=\sum_{i=1}^m\left(\sum_{j=1}^n a_{ij}x_j\right)w_i.
$$
따라서 $[T(x)]_{\mathcal C}$의 $i$번째 성분은 $\sum_{j=1}^n a_{ij}x_j$이고, 이것은 정확히 $A[x]_{\mathcal B}$의 $i$번째 성분이다. 그러므로
$$
[T(x)]_{\mathcal C}=A[x]_{\mathcal B}.
$$
\end{proof}

\paragraph{의미 해석.}
이 정리는 ``선형사상 계산''을 ``행렬곱 계산''으로 바꾸는 핵심 규칙입니다. 이후에는 사상을 직접 다루기보다, 적절한 기저를 정해 행렬로 옮겨 계산할 수 있습니다.

\begin{theorem}[합성사상의 행렬공식]
선형사상 $T:V\to W$, $S:W\to U$와 순서기저 $\mathcal B,\mathcal C,\mathcal D$에 대해
$$
[S\circ T]_{\mathcal D\leftarrow\mathcal B}
=
[S]_{\mathcal D\leftarrow\mathcal C}
[T]_{\mathcal C\leftarrow\mathcal B}
$$
가 성립한다.
\end{theorem}

\paragraph{증명 전략.}
임의의 $x\in V$에 대해 좌표식을 두 번 적용합니다. 한편 $(S\circ T)$에 대한 좌표식도 적용해 같은 벡터를 두 방식으로 표현하고, 좌표사상의 전사성을 이용해 행렬을 동일시합니다.

\begin{proof}
임의의 $x\in V$를 잡는다. 좌표식을 $T$와 $S$에 차례로 적용하면
$$
[(S\circ T)(x)]_{\mathcal D}
=[S(T(x))]_{\mathcal D}
=[S]_{\mathcal D\leftarrow\mathcal C}[T(x)]_{\mathcal C}
=[S]_{\mathcal D\leftarrow\mathcal C}[T]_{\mathcal C\leftarrow\mathcal B}[x]_{\mathcal B}.
$$
또한 $(S\circ T):V\to U$에 좌표식을 적용하면
$$
[(S\circ T)(x)]_{\mathcal D}
=
[S\circ T]_{\mathcal D\leftarrow\mathcal B}[x]_{\mathcal B}.
$$
따라서 모든 $x\in V$에 대해
$$
[S\circ T]_{\mathcal D\leftarrow\mathcal B}[x]_{\mathcal B}
=
[S]_{\mathcal D\leftarrow\mathcal C}[T]_{\mathcal C\leftarrow\mathcal B}[x]_{\mathcal B}.
$$
좌표사상 $x\mapsto [x]_{\mathcal B}$는 $\mathbb F^n$으로의 전단사이므로 우변과 좌변은 모든 $y\in\mathbb F^n$에 대해 같은 값을 갖는다. 따라서 두 행렬이 같아
$$
[S\circ T]_{\mathcal D\leftarrow\mathcal B}
=
[S]_{\mathcal D\leftarrow\mathcal C}
[T]_{\mathcal C\leftarrow\mathcal B}.
$$
\end{proof}

\paragraph{의미 해석.}
행렬곱의 순서가 합성의 순서와 정확히 대응합니다. 그래서 복잡한 사상의 연쇄작용도 적절한 중간 기저를 선택하면 체계적으로 계산할 수 있습니다.

\begin{example}
$T:\mathbb R^2\to\mathbb R^2$, $T(x,y)=(2x-y,\;x+3y)$를 두고
$$
\mathcal B=((1,1),(1,-1)),\qquad \mathcal C=((1,0),(0,1))
$$
라 하자. 그러면
$$
T(1,1)=(1,4),\qquad T(1,-1)=(3,-2)
$$
이므로
$$
[T]_{\mathcal C\leftarrow\mathcal B}
=
\begin{bmatrix}
1&3\\
4&-2
\end{bmatrix}.
$$
즉 $\mathcal B$-좌표 $(a,b)^T$를 가지는 벡터는 $T$에 의해 $\mathcal C$-좌표
$\begin{bmatrix}1&3\\4&-2\end{bmatrix}\begin{bmatrix}a\\b\end{bmatrix}$
로 이동한다.
\end{example}

\subsection*{주의}
\begin{warning}
표현행렬의 열은 ``기저벡터의 영상 좌표''입니다. 행을 기저영상으로 읽으면 정의와 어긋납니다.
\end{warning}
\begin{warning}
$[T]_{\mathcal C\leftarrow\mathcal B}$에서 화살표 방향은 필수 정보입니다. 정의역 기저와 공역 기저를 바꾸어 쓰면 같은 숫자 배열이라도 전혀 다른 의미가 됩니다.
\end{warning}

\subsection*{자가진단퀴즈}
\begin{enumerate}[label=\arabic*.]
\item $T:\mathbb R^2\to\mathbb R^2$, $T(x,y)=(x+2y,3x)$, $\mathcal B=((1,0),(1,1))$, $\mathcal C$ 표준기저일 때 $[T]_{\mathcal C\leftarrow\mathcal B}$를 구해 보십시오.
\item 좌표식 정리의 증명에서 ``좌표사상이 전단사''라는 사실이 어디에서 사용되는지 설명해 보십시오.
\item $R:\mathbb R^2\to\mathbb R^2$, $S:\mathbb R^2\to\mathbb R^2$를 임의로 잡고 $[R\circ S]$의 곱셈 순서가 왜 $[R][S]$인지 말로 설명해 보십시오.
\end{enumerate}

\sectiontemplate{기저변환(change of basis)과 유사공식}{같은 공간의 서로 다른 기저는 같은 벡터를 다른 좌표로 기록합니다. 이 좌표변환이 선형사상 행렬에 어떤 변화를 주는지 정리하면 유사공식이 자연스럽게 나옵니다.}{\item 기저변환행렬의 정의와 가역성을 설명할 수 있습니다.\item 좌표변환식 $[x]_{\mathcal B}=P[x]_{\mathcal B'}$를 증명할 수 있습니다.\item 같은 연산자의 두 표현행렬이 $A'=P^{-1}AP$를 만족함을 유도할 수 있습니다.}{\item 역행렬\item 항등사상의 행렬표현}

\begin{definition}
$V$의 두 순서기저 $\mathcal B,\mathcal B'$에 대해
$$
P_{\mathcal B\leftarrow\mathcal B'}
:=
[I_V]_{\mathcal B\leftarrow\mathcal B'}
$$
를 $\mathcal B'$-좌표를 $\mathcal B$-좌표로 바꾸는 기저변환행렬이라 한다.
\end{definition}

\begin{theorem}[기저변환행렬의 기본 성질]
$P=P_{\mathcal B\leftarrow\mathcal B'}$, $Q=P_{\mathcal B'\leftarrow\mathcal B}$라 하면
$$
PQ=I,\qquad QP=I
$$
이므로 $P$는 가역이고 $P^{-1}=Q$이다. 또한 모든 $x\in V$에 대해
$$
[x]_{\mathcal B}=P[x]_{\mathcal B'}
$$
가 성립한다.
\end{theorem}

\paragraph{증명 전략.}
항등사상 $I_V$의 표현행렬을 합성공식에 대입해 $PQ=I$, $QP=I$를 얻습니다. 좌표변환식은 $I_V$에 대한 좌표식 정리를 그대로 적용하면 바로 나옵니다.

\begin{proof}
합성공식을 $I_V=I_V\circ I_V$에 적용하면
$$
[I_V]_{\mathcal B\leftarrow\mathcal B}
=
[I_V]_{\mathcal B\leftarrow\mathcal B'}
[I_V]_{\mathcal B'\leftarrow\mathcal B}
$$
이므로
$$
I=PQ.
$$
같은 방식으로
$$
[I_V]_{\mathcal B'\leftarrow\mathcal B'}
=
[I_V]_{\mathcal B'\leftarrow\mathcal B}
[I_V]_{\mathcal B\leftarrow\mathcal B'}
$$
에서
$$
I=QP
$$
를 얻는다. 따라서 $P$는 가역이고 $P^{-1}=Q$이다.

이제 임의의 $x\in V$에 대해 좌표식 정리를 $I_V$에 적용하면
$$
[x]_{\mathcal B}
=
[I_V]_{\mathcal B\leftarrow\mathcal B'}[x]_{\mathcal B'}
=
P[x]_{\mathcal B'}
$$
가 성립한다.
\end{proof}

\paragraph{의미 해석.}
기저변환행렬은 단순한 계산도구가 아니라 ``같은 벡터의 다른 기록법'' 사이를 연결하는 사상입니다. 가역성은 정보 손실 없이 좌표를 바꿀 수 있음을 뜻합니다.

\begin{theorem}[같은 연산자의 두 표현행렬]
$T\in\operatorname{End}(V)$라 하고
$$
A=[T]_{\mathcal B},\qquad A'=[T]_{\mathcal B'},\qquad
P=P_{\mathcal B\leftarrow\mathcal B'}
$$
라 하자. 그러면
$$
A'=P^{-1}AP
$$
가 성립한다.
\end{theorem}

\paragraph{증명 전략.}
임의의 벡터 $x$에 대해 $[x]_{\mathcal B}=P[x]_{\mathcal B'}$와 $[T(x)]_{\mathcal B}=P[T(x)]_{\mathcal B'}$를 동시에 사용합니다. 같은 $[T(x)]_{\mathcal B}$를 두 식으로 표현해 $PA'=AP$를 얻고, 마지막에 역행렬을 곱합니다.

\begin{proof}
임의의 $x\in V$를 잡는다. 좌표변환식으로
$$
[x]_{\mathcal B}=P[x]_{\mathcal B'}
$$
이고, $T(x)$에 대해서도
$$
[T(x)]_{\mathcal B}=P[T(x)]_{\mathcal B'}
$$
이다. 한편 정의에 의해
$$
[T(x)]_{\mathcal B}=A[x]_{\mathcal B},\qquad
[T(x)]_{\mathcal B'}=A'[x]_{\mathcal B'}.
$$
따라서
$$
A[x]_{\mathcal B}=P A'[x]_{\mathcal B'}.
$$
여기에 $[x]_{\mathcal B}=P[x]_{\mathcal B'}$를 대입하면
$$
AP[x]_{\mathcal B'}=PA'[x]_{\mathcal B'}
$$
가 모든 $x\in V$에 대해 성립한다. 좌표사상의 전사성을 이용하면
$$
AP=PA'
$$
이고, 양변 왼쪽에 $P^{-1}$를 곱해
$$
A'=P^{-1}AP
$$
를 얻는다.
\end{proof}

\paragraph{의미 해석.}
유사공식은 ``기저만 바꿨을 때 행렬이 어떻게 변하는가''를 완전히 기술합니다. 즉 숫자 배열의 변화는 본질적 변화가 아니라 좌표계 변경 효과입니다.

\begin{example}
$T(x,y)=(2x-y,\;x+3y)$의 표준기저 표현행렬은
$$
A=\begin{bmatrix}2&-1\\1&3\end{bmatrix}
$$
이다. $\mathcal B'=((1,1),(1,-1))$에 대해
$$
P=P_{\mathrm{std}\leftarrow\mathcal B'}=
\begin{bmatrix}1&1\\1&-1\end{bmatrix},\qquad
P^{-1}=\frac12\begin{bmatrix}1&1\\1&-1\end{bmatrix}
$$
이므로
$$
A'=P^{-1}AP=
\begin{bmatrix}
\frac52&\frac12\\[2mm]
-\frac32&\frac52
\end{bmatrix}.
$$
이 $A'$는 같은 연산자 $T$를 $\mathcal B'$에서 본 표현입니다.
\end{example}

\subsection*{주의}
\begin{warning}
$P_{\mathcal B\leftarrow\mathcal B'}$와 $P_{\mathcal B'\leftarrow\mathcal B}$를 혼동하면 유사공식의 역행렬 위치가 뒤바뀝니다. 화살표 방향을 먼저 고정한 뒤 계산하십시오.
\end{warning}
\begin{warning}
$A'=P^{-1}AP$에서 $P$는 ``새 기저벡터를 옛 기저로 적은 열벡터''를 모아 만든 행렬입니다. 임의의 가역행렬을 대입하면 기저 해석이 깨집니다.
\end{warning}

\subsection*{자가진단퀴즈}
\begin{enumerate}[label=\arabic*.]
\item 두 기저 $\mathcal B,\mathcal B'$가 주어졌을 때 $P_{\mathcal B\leftarrow\mathcal B'}$의 각 열이 무엇을 뜻하는지 설명해 보십시오.
\item $A' = P^{-1}AP$를 $AP=PA'$ 형태로 다시 쓰면 어떤 계산상 이점이 있는지 말해 보십시오.
\item $T$의 두 표현행렬 $A,A'$가 주어졌을 때, 어떤 정보를 확인하면 ``같은 연산자의 다른 표현''인지 판정할 수 있는지 정리해 보십시오.
\end{enumerate}

\sectiontemplate{유사(similarity) 관계와 동치류}{유사는 ``기저변환으로 서로 바뀌는가''를 행렬만으로 표현한 관계입니다. 이 절에서는 유사가 동치관계를 이루며, 선형연산자 분류 문제를 동치류 분류 문제로 바꿔 준다는 점을 확인하겠습니다.}{\item 유사의 정의를 정확히 진술할 수 있습니다.\item 유사가 동치관계임을 완전하게 증명할 수 있습니다.\item 유사가 ``같은 연산자의 서로 다른 표현''과 동치임을 설명할 수 있습니다.}{\item 가역행렬과 역행렬\item 동치관계의 세 성질}

\begin{definition}
$A,B\in M_n(\mathbb F)$에 대해 어떤 가역행렬 $P\in GL_n(\mathbb F)$가 존재하여
$$
B=P^{-1}AP
$$
가 성립하면 $A$와 $B$가 \emph{유사(similar)}라고 하고, $A\sim B$로 쓴다.
\end{definition}

\begin{theorem}[유사는 동치관계]
$M_n(\mathbb F)$ 위의 관계 $\sim$은 반사성, 대칭성, 추이성을 모두 만족한다.
\end{theorem}

\paragraph{증명 전략.}
각 성질을 정의식 $B=P^{-1}AP$에 직접 대입해 확인합니다. 반사성은 $P=I$, 대칭성은 역행렬 사용, 추이성은 두 식을 연쇄 대입하고 $(PQ)^{-1}=Q^{-1}P^{-1}$을 사용합니다.

\begin{proof}
\textbf{(반사성)} 임의의 $A$에 대해
$$
A=I^{-1}AI
$$
이므로 $A\sim A$.

\textbf{(대칭성)} $A\sim B$라 하여 $B=P^{-1}AP$라 하자. 그러면
$$
A=PBP^{-1}=(P^{-1})^{-1}B(P^{-1})
$$
이므로 $B\sim A$.

\textbf{(추이성)} $A\sim B$, $B\sim C$라 하자. 그러면 어떤 가역행렬 $P,Q$에 대해
$$
B=P^{-1}AP,\qquad C=Q^{-1}BQ.
$$
대입하면
$$
C=Q^{-1}(P^{-1}AP)Q=(PQ)^{-1}A(PQ).
$$
$PQ$도 가역이므로 $A\sim C$.

세 성질이 모두 성립하므로 $\sim$은 동치관계이다.
\end{proof}

\paragraph{의미 해석.}
유사는 행렬 전체를 ``같은 본질을 가진 묶음''으로 분할합니다. 이후의 정준형 이론은 각 동치류에서 대표를 찾는 작업으로 이해할 수 있습니다.

\begin{theorem}[유사의 연산자 해석]
$A,B\in M_n(\mathbb F)$에 대해 다음이 동치이다.
\begin{enumerate}[label=(\arabic*)]
\item $A\sim B$.
\item 어떤 $n$차원 벡터공간 $V$, 선형연산자 $T\in\operatorname{End}(V)$, 두 순서기저 $\mathcal E,\mathcal F$가 존재하여
$$
A=[T]_{\mathcal E},\qquad B=[T]_{\mathcal F}
$$
가 성립한다.
\end{enumerate}
\end{theorem}

\paragraph{증명 전략.}
$(2)\Rightarrow(1)$은 기저변환 공식으로 즉시 보입니다. $(1)\Rightarrow(2)$에서는 $V=\mathbb F^n$ 위에 $T(x)=Ax$를 정의하고, 유사식의 가역행렬 $P$의 열벡터를 새 기저로 잡아 $B$가 새 기저 표현행렬임을 보입니다.

\begin{proof}
\textbf{$(2)\Rightarrow(1)$:}
가정에서 $A=[T]_{\mathcal E}$, $B=[T]_{\mathcal F}$이다.
$P=P_{\mathcal E\leftarrow\mathcal F}$라 두면 앞 절의 유사공식에 의해
$$
B=P^{-1}AP
$$
가 되어 $A\sim B$.

\textbf{$(1)\Rightarrow(2)$:}
$A\sim B$라 하여 어떤 가역행렬 $P$에 대해 $B=P^{-1}AP$라 하자.
$V=\mathbb F^n$로 두고 $T:V\to V$를 $T(x)=Ax$로 정의한다.
표준기저를 $\mathcal E$라 하면 $[T]_{\mathcal E}=A$이다.

이제 $P$의 열벡터를 $u_1,\dots,u_n$이라 두고 $\mathcal F=(u_1,\dots,u_n)$라 하자.
$P$가 가역이므로 열벡터들은 일차독립이며 $\mathcal F$는 $V$의 기저이다.
또한 정의에 의해
$$
P=P_{\mathcal E\leftarrow\mathcal F}.
$$
따라서 유사공식을 적용하면
$$
[T]_{\mathcal F}=P^{-1}[T]_{\mathcal E}P=P^{-1}AP=B.
$$
즉 (2)가 성립한다.
\end{proof}

\paragraph{의미 해석.}
유사는 임의의 대수적 변환이 아니라 ``같은 연산자를 다른 좌표계에서 본 결과''와 정확히 일치합니다. 이 관점은 유사 불변량을 이해하는 기준점이 됩니다.

\begin{theorem}[항등행렬과 영행렬의 고정성]
$A\sim I_n$이면 $A=I_n$이고, $A\sim 0_n$이면 $A=0_n$이다.
\end{theorem}

\paragraph{증명 전략.}
유사식에 직접 대입합니다. $I_n$과 $0_n$은 가역행렬과의 공액변환에서 형태가 바뀌지 않습니다.

\begin{proof}
$A\sim I_n$이면 어떤 가역행렬 $P$에 대해
$$
A=P^{-1}I_nP=I_n
$$
이다.
또 $A\sim 0_n$이면
$$
A=P^{-1}0_nP=0_n
$$
이다.
\end{proof}

\paragraph{의미 해석.}
모든 행렬이 어떤 표준형으로 단순화되는 것은 아닙니다. 특히 항등행렬과 영행렬은 유사동치류가 한 원소만 있는 극단적 사례입니다.

\begin{example}
$$
A=\begin{bmatrix}1&0\\0&2\end{bmatrix},\qquad
P=\begin{bmatrix}1&1\\0&1\end{bmatrix}
$$
라 하면
$$
P^{-1}=\begin{bmatrix}1&-1\\0&1\end{bmatrix},\qquad
B=P^{-1}AP=\begin{bmatrix}1&-1\\0&2\end{bmatrix}.
$$
따라서 $A\sim B$이다. 즉 대각행렬이 아닌 형태라도 같은 연산자의 다른 좌표표현일 수 있다.
\end{example}

\subsection*{주의}
\begin{warning}
유사는 같은 크기의 정사각행렬 사이에서만 정의됩니다. $m\times n$ 행렬($m\neq n$)에는 유사를 적용할 수 없습니다.
\end{warning}
\begin{warning}
행동치(row equivalence)와 유사(similarity)는 전혀 다른 관계입니다. 행동치는 연립방정식 해집합과 연결되고, 유사는 선형연산자의 기저변환과 연결됩니다.
\end{warning}

\subsection*{자가진단퀴즈}
\begin{enumerate}[label=\arabic*.]
\item 유사의 동치관계 증명에서 추이성 단계에 왜 $(PQ)^{-1}=Q^{-1}P^{-1}$이 필요한지 설명해 보십시오.
\item ``$A\sim B$''를 ``같은 선형연산자의 두 표현''으로 해석하는 과정을 $V=\mathbb F^n$에서 직접 써 보십시오.
\item $A\sim I_n$이면 $A=I_n$임을 유사식 한 줄로 증명해 보십시오.
\end{enumerate}

\sectiontemplate{유사 불변량(invariant)과 판정 도구}{유사동치류를 구분하려면 바뀌지 않는 양이 필요합니다. 이 절에서는 특성다항식, 행렬식, trace, rank 계열과 같은 기본 불변량을 엄밀하게 확인하고, 무엇이 충분조건이 아닌지도 함께 점검하겠습니다.}{\item 유사 불변량의 정의를 쓸 수 있습니다.\item 특성다항식과 그 계수들이 유사에서 보존됨을 증명할 수 있습니다.\item $\,\operatorname{rank}(A-\lambda I)^k$ 계열을 사용해 비유사를 판정할 수 있습니다.}{\item 행렬식의 곱셈성\item rank-nullity 정리}

\begin{definition}
함수 $f:M_n(\mathbb F)\to S$가
$$
A\sim B \Longrightarrow f(A)=f(B)
$$
를 만족하면 $f$를 유사에 대한 \emph{불변량(invariant)}이라 한다.
\end{definition}

\begin{theorem}[특성다항식 불변성]
$A,B\in M_n(\mathbb F)$가 유사하면
$$
\chi_A(t)=\chi_B(t)
$$
가 성립한다. 따라서 $\det A=\det B$, $\operatorname{tr}A=\operatorname{tr}B$이다.
\end{theorem}

\paragraph{증명 전략.}
$B=P^{-1}AP$를 가정하고 $tI-B$를 $P^{-1}(tI-A)P$ 형태로 바꿉니다. 행렬식의 곱셈성을 적용하면 특성다항식이 같아집니다. trace와 determinant는 특성다항식 계수로부터 얻습니다.

\begin{proof}
$B=P^{-1}AP$라 하자. 그러면
$$
tI-B=tI-P^{-1}AP=P^{-1}(tI-A)P
$$
이다. 따라서
$$
\chi_B(t)=\det(tI-B)
=\det\!\big(P^{-1}(tI-A)P\big)
=\det(P^{-1})\det(tI-A)\det(P)
=\det(tI-A)
=\chi_A(t).
$$
즉 $\chi_A(t)=\chi_B(t)$이다.

또한
$$
\chi_A(t)=t^n-(\operatorname{tr}A)t^{n-1}+\cdots+(-1)^n\det A,
$$
$$
\chi_B(t)=t^n-(\operatorname{tr}B)t^{n-1}+\cdots+(-1)^n\det B
$$
이므로 두 다항식이 같다는 사실에서 계수 비교를 하면
$$
\operatorname{tr}A=\operatorname{tr}B,\qquad \det A=\det B
$$
를 얻는다.
\end{proof}

\paragraph{의미 해석.}
고유값 구조를 담는 특성다항식이 보존되므로, 유사는 연산자의 스펙트럼 정보를 바꾸지 않습니다. 다만 같은 특성다항식이 항상 유사를 보장하지는 않습니다.

\begin{theorem}[rank 계열 불변성]
$B=P^{-1}AP$라 하자. 임의의 $\lambda\in\mathbb F$와 $k\in\mathbb N$에 대해
$$
(B-\lambda I)^k=P^{-1}(A-\lambda I)^kP
$$
이며, 따라서
$$
\operatorname{rank}(B-\lambda I)^k=\operatorname{rank}(A-\lambda I)^k
$$
및
$$
\operatorname{nullity}(B-\lambda I)^k=\operatorname{nullity}(A-\lambda I)^k
$$
가 성립한다.
\end{theorem}

\paragraph{증명 전략.}
먼저 $B-\lambda I=P^{-1}(A-\lambda I)P$를 확인하고, 거듭제곱은 귀납으로 처리합니다. 그 다음 상(image)의 정확한 관계를 써서 rank 보존을 보이고, nullity는 rank-nullity로 결론냅니다.

\begin{proof}
$B=P^{-1}AP$에서
$$
B-\lambda I=P^{-1}AP-\lambda I=P^{-1}(A-\lambda I)P
$$
이다. 이에 대해 귀납법으로
$$
(B-\lambda I)^k=P^{-1}(A-\lambda I)^kP
$$
를 얻는다.

이제 상공간을 비교한다. 임의의 $y\in\operatorname{im}(B-\lambda I)^k$에 대해 어떤 $x$가 존재하여
$$
y=(B-\lambda I)^k x=P^{-1}(A-\lambda I)^k(Px),
$$
즉 $y\in P^{-1}\!\big(\operatorname{im}(A-\lambda I)^k\big)$이다. 반대로 임의의
$y=P^{-1}(A-\lambda I)^k z$에 대해 $x=P^{-1}z$를 택하면
$$
y=(B-\lambda I)^k x
$$
이므로
$$
\operatorname{im}(B-\lambda I)^k
=
P^{-1}\!\big(\operatorname{im}(A-\lambda I)^k\big).
$$
$P^{-1}$는 동형이므로 두 상공간의 차원이 같아
$$
\operatorname{rank}(B-\lambda I)^k=\operatorname{rank}(A-\lambda I)^k
$$
이다.

두 연산자 $(B-\lambda I)^k,(A-\lambda I)^k$는 모두 $\mathbb F^n\to\mathbb F^n$이므로 rank-nullity 정리에 의해
$$
\operatorname{nullity}(X)=n-\operatorname{rank}(X)
$$
가 성립한다. 따라서 nullity도 서로 같다.
\end{proof}

\paragraph{의미 해석.}
$\operatorname{rank}(A-\lambda I)^k$ 계열은 단순한 $\det,\operatorname{tr}$보다 훨씬 세밀한 정보입니다. 이후 Jordan 이론에서 블록 크기를 판정하는 핵심 지표로 쓰입니다.

\begin{example}
다음 두 행렬을 비교하자.
$$
A=I_2,\qquad
C=\begin{bmatrix}1&1\\0&1\end{bmatrix}.
$$
둘 다 $\chi(t)=(t-1)^2$, $\det=1$, $\operatorname{tr}=2$를 가진다. 그러나
$$
\operatorname{rank}(A-I_2)=0,\qquad
\operatorname{rank}(C-I_2)=1
$$
이므로 유사일 수 없다. 즉 특성다항식만으로는 유사를 판정할 수 없다.
\end{example}

\subsection*{주의}
\begin{warning}
$\det$와 $\operatorname{tr}$가 같다는 사실은 필요조건일 뿐 충분조건이 아닙니다. 이 둘만으로 유사 여부를 결론내리면 자주 오답이 납니다.
\end{warning}
\begin{warning}
특성다항식이 같아도 Jordan 구조가 다르면 유사하지 않습니다. 실제 판정에서는 $\operatorname{rank}(A-\lambda I)^k$ 같은 추가 불변량이 필요합니다.
\end{warning}

\subsection*{자가진단퀴즈}
\begin{enumerate}[label=\arabic*.]
\item $B=P^{-1}AP$에서 $\chi_B(t)=\chi_A(t)$를 행렬식 곱셈성만 사용해 직접 유도해 보십시오.
\item 왜 $\operatorname{rank}(A-\lambda I)$가 유사 불변량인지 한 줄 식으로 설명해 보십시오.
\item $\chi_A=\chi_B$이지만 $A\not\sim B$인 $2\times2$ 예를 하나 만들고, 비유사 근거를 제시해 보십시오.
\end{enumerate}

\chapterapplicationtemplate{선형 미분방정식 $x'(t)=Ax(t)$에서 좌표변환 $x=Py$를 적용하면
$$
y'(t)=P^{-1}AP\,y(t)
$$
를 얻습니다. 따라서 해의 질적 거동(고유값, 지수성장률)은 좌표계와 무관하며, 계산 편의를 위해 가장 다루기 쉬운 기저를 선택할 수 있습니다. 실제로 한 시스템을 두 기저에서 행렬로 표현해 보고, 유사 불변량이 어떻게 동일하게 유지되는지 확인해 보십시오.}

\chaptersummarytemplate

\section*{종합 연습문제}

\subsection*{연습문제 A (기초 확인)}
\begin{enumerate}[label=A\arabic*.]
\item $T:\mathbb R^2\to\mathbb R^2$, $T(x,y)=(3x+y,2x-y)$, $\mathcal B=((1,0),(1,1))$, $\mathcal C$ 표준기저일 때 $[T]_{\mathcal C\leftarrow\mathcal B}$를 구하시오.
\item $\mathcal B=((1,1),(1,-1))$에서 좌표가 $\begin{bmatrix}2\\-1\end{bmatrix}$인 벡터를 표준좌표로 바꾸시오.
\item $\mathcal B=((1,0),(0,1))$, $\mathcal B'=((1,1),(1,-1))$에 대해 $P_{\mathcal B\leftarrow\mathcal B'}$와 $P_{\mathcal B'\leftarrow\mathcal B}$를 구하시오.
\item $A=\begin{bmatrix}2&-1\\1&3\end{bmatrix}$, $P=\begin{bmatrix}1&1\\1&-1\end{bmatrix}$일 때 $P^{-1}AP$를 계산하시오.
\item $A\sim B$이면 $\det A=\det B$임을 유사식 한 줄로 확인하시오.
\item $A\sim B$이면 $\operatorname{tr}A=\operatorname{tr}B$임을 특성다항식 계수 비교로 설명하시오.
\item $A=\begin{bmatrix}1&0\\0&1\end{bmatrix}$, $B=\begin{bmatrix}1&1\\0&1\end{bmatrix}$에 대해 $\operatorname{rank}(A-I)$, $\operatorname{rank}(B-I)$를 구하고 유사 여부를 판정하시오.
\item $A\sim 0_3$이면 $A=0_3$임을 증명하시오.
\end{enumerate}

\subsection*{연습문제 B (개념 연결)}
\begin{enumerate}[label=B\arabic*.]
\item $T:V\to W$의 표현행렬이 $A=[T]_{\mathcal C\leftarrow\mathcal B}$일 때, 기저를 $\mathcal B',\mathcal C'$로 바꾸면
$$
[T]_{\mathcal C'\leftarrow\mathcal B'}=
P_{\mathcal C'\leftarrow\mathcal C}\,A\,P_{\mathcal B\leftarrow\mathcal B'}
$$
가 됨을 증명하시오.
\item 같은 특성다항식을 가지지만 유사하지 않은 $2\times2$ 행렬의 예를 하나 제시하고, 왜 유사하지 않은지 불변량으로 설명하시오.
\item $A\in M_n(\mathbb F)$와 가역행렬 $P$에 대해 $A$와 $P^{-1}AP$의 고유값(대수적 중복도 포함)이 일치함을 보이시오.
\item $A$가 대각행렬 $D$와 유사하다고 하자. $P^{-1}AP=D$에서 $P$의 열벡터가 $A$의 고유벡터임을 증명하시오.
\item $N=\begin{bmatrix}0&1\\0&0\end{bmatrix}$, $0_2$는 둘 다 $\det=0$, $\operatorname{tr}=0$이다. 두 행렬이 유사하지 않음을 $\operatorname{rank}(X)$ 또는 $\operatorname{rank}(X-I)$ 중 적절한 불변량으로 판정하시오.
\end{enumerate}

\subsection*{연습문제 C (증명/종합)}
\begin{enumerate}[label=C\arabic*.]
\item $B=P^{-1}AP$이고 $p(t)\in\mathbb F[t]$일 때
$$
p(B)=P^{-1}p(A)P
$$
임을 증명하시오.
\item ``대각화 가능성은 유사 불변량''임을 증명하시오. 즉 $A\sim B$이면 $A$가 대각화 가능할 필요충분조건은 $B$가 대각화 가능한 것이다.
\item $A,B\in M_3(\mathbb C)$가 같은 특성다항식 $(t-\lambda)^3$를 가진다고 하자. $\operatorname{rank}(A-\lambda I)$, $\operatorname{rank}(A-\lambda I)^2$와 같은 수열이 Jordan 블록 크기를 결정한다는 사실을 이용해, 유사 판정 절차를 서술하시오.
\end{enumerate}